% !TEX program = xelatex
\documentclass[
    school=etsisi,
    type=pfg,
    degree=61CI,
]{upm-report}
%%%%%%%%%%%%%%%%%%%%%%%%%%%%%%%%%%%%%%%%%%%%%%%%%%%%%%%%%%%%%%%%%%%%%%%%
% TÍTULO, AUTORES Y DIRECTORES
%
\title{Plantilla de proyecto de fin de estudios para la UPM}
\author{Pedro Mario Cea Torralba}
\bibauthor{}
\director{}
% TODO: añadir el director

%%%%%%%%%%%%%%%%%%%%%%%%%%%%%%%%%%%%%%%%%%%%%%%%%%%%%%%%%%%%%%%%%%%%%%%%
% RESUMEN Y ABSTRACT
%
\abstract{spanish}{
    El resumen de un \acrlong{pfg} o de un \acrlong{pfm} condensa en
    tres o cuatro párrafos el contenido de la memoria. Se debería asumir
    que el lector tiene una cierta idea de lo que se trata y también que
    si no le enganchamos, probablemente pase de él (eso es malo).
    
    \textbf{Condensado no quiere decir incompleto}. Debe contener la
    información más destacable. Lo ideal: que ocupe entre media y una
    cara de un folio A4. Comenzará por el propósito y principales
    objetivos de la memoria, seguirá con los aspectos más destacables
    de la metodología empleada, y luego de los resultados obtenidos. Por
    último se presentarán la o las conclusiones más importantes a las
    que se ha llegado.
    
    Que tenga un estilo claro y conciso, sin ambigüedades de ningún tipo.
    Además, al ser un resumen de todo el contenido, ni que decir tiene
    que deberá ser lo último que escribamos, y deberá mantener una
    absoluta fidelidad con el contenido de la memoria.
}
\keywords{spanish}{
    Cuatro o cinco;
    Expresiones clave;
    Que clasifiquen;
    Nuestro proyecto o;
    Investigación
}

\abstract{english}{
    This section must contain the summary that we have written before in
    Spanish, but in English, as well as the keywords.
}
\keywords{english}{
    Four or five;
    Key Expressions;
    Summarising;
    Our Project or;
    Research
}

% Opcional
\reportquotation{
    Oigo y olvido, veo y recuerdo, hago y aprendo.
}{Confucio}

% Opcional
\acknowledgements{
    Agradecimientos que se quieran dar. Si no se quiere añadir
    agradecimientos, es tan sencillo como eliminar la entrada
    \texttt{\textbackslash acknowledgements} del fichero
    \texttt{report.tex}.
}

%%%%%%%%%%%%%%%%%%%%%%%%%%%%%%%%%%%%%%%%%%%%%%%%%%%%%%%%%%%%%%%%%%%%%%%%
% GLOSARIO Y ABREVIATURAS
%
\newglossaryentry{arco-perdicion}{
        name=Arco de la Perdición,
        description={Arco único que posee poderes destructivos capaz de
            desatar grandes catástrofes y de traer desgracias a aquellos
            que se encuentren en su camino
        }
}
\newglossaryentry{hacha-batalla}{
        name=hacha de batalla,
        description={Herramienta antigua utilizada en combate,
            caracterizada por su doble función de arma y herramienta
        }
}
\newglossaryentry{python}{
    name={Python},
    plural={Pythonacos},
    description={El mejor lenguaje de programación}
}

\newabbreviation[
    longplural={inteligencias artificiales}
    ]{ia}{IA}{inteligencia artificial}
\newabbreviation[
    longplural={organizaciones no gubernamentales},
    shortplural={ONG}
    ]{ong}{ONG}{organización no gubernamental}
\newabbreviation[
    longplural={proyectos de fin de grado}
    shortplural={PFG}
    ]{pfg}{PFG}{proyecto de fin de grado}
\newabbreviation[
    longplural={proyectos de fin de máster}
    shortplural={PFM}
    ]{pfm}{PFM}{proyecto de fin de máster}
\newabbreviation[
    description={Del inglés \textit{role-playing game}},
    shortplural={RPG}
    ]{rpg}{RPG}{juego de rol}
\newabbreviation[
    description={S.P.E.C.I.A.L es la sigla usada para los atributos de
        Fuerza (\textbf{S}trenght), Percepción (\textbf{P}erception),
        Resistencia (\textbf{E}ndurance), Carisma, (\textbf{C}harisma),
        Inteligencia, (\textbf{I}ntelligence), Agilidad,
        (\textbf{A}gility), y Suerte (\textbf{L}uck)
    },
    ]{special}{SPECIAL}{\textit{Strenght, Perception, Endurance,
        Charisma, Intelligence, Agility \& Luck}
    }

\begin{document}

\chapter{Introducción}
\label{ch:introduccion}

La introducción a un trabajo de fin de estudios es el punto de entrada a
todo el trabajo realizado y es considerada la más importante tras el
resumen, donde se resume el trabajo entero. Hay que dejar claro
\textbf{qué} se ha hecho y el \textbf{porqué} de su importancia. Intenta
generar expectativa. Un gancho típico suele ser discutir sobre un dato
relevante o controvertido o plantear una pregunta relevante en lo que se
está trabajando.

\section{Objetivos}

Esta sección es \textbf{obligatoria}, y determinas la \textbf{finalidad}
del proyecto. Suele encajar en una de las siguientes categorías:

\begin{itemize}
    \item \textbf{Contraste} o validación de una hipótesis, típico de en
        \glspl{pfm},
    \item \textbf{Desarrollo} o diseño de algo (\textit{hardware},
        sistemas, edificios), común en ingenierías,
    \item \textbf{Estudio} de un tema que deduce o descubre nuevo
        conocimiento, típico en ciencias puras y humanidades.
\end{itemize}

Nos vienen muy bien para evaluar el éxito del proyecto, porque así en
las conclusiones podemos indicar qué objetivos se han logrado y cuáles
no, y el porqué. Para poder decir si un objetivo se ha cumplido ha de
ser (i) \textbf{acotado en el tiempo} ---así se puede planificar---,
(ii) \textbf{medible} ---podemos cuantificar su éxito---, (iii)
\textbf{específico} ---evitando así tareas irrelevantes---, (iv)
\textbf{Alcanzable}  ---asegurando su finalización--- y (v)
\textbf{relevante} ---es decir, relacionado con el campo de estudio---.

\section{Motivación}
\label{s:motivacion}

Decidí decantarme por este proyecto, ya que junta un \textbf{problema real} que existe en nuestra sociedad. Hay gente que no se conforma a los teclados clásicos, ya que su cuerpo no puede escribir cómodamente con ellos. Estas personas tienen que soportar esa situación e intentar adaptarse a herramientas que tendrían que adaptarse a ellas.

Es cierto que existen alternativas, como poder hablar y que se transcriba el texto. Sin embargo, no es una solución suficientemente buena para todos los casos. Por ejemplo, Si un estudiante está en clase tomando apuntes, mientras que el profesor está explicando, este no podría, ya que estaría interfiriendo con la clase. En estas situaciones no se puede estar hablando.

Entonces, eso hace que tengas que estar usando métodos en los que no vas a poder escribir de forma tan eficiente. Incluso podrían llegar a ser métodos \textbf{ineficientes, incómodos o hasta lesivos}.

Por otro lado, mi motivación para hacer este proyecto ha sido que, mientras avanzaba en la carrera, me he dado cuenta de que este aspecto de los \textbf{sistemas embarcados o empotrados} es lo que más me ha interesado y más me ha gustado.

Creo que la posibilidad de juntar lo que me apasiona con el poder ayudar a otros y que puedan vivir una vida mejor es algo que encuentro realmente bonito.

\section{Justificación}
\label{sec:justificacion}

La justificación de este trabajo se basa en la necesidad de explorar alternativas de interacción que puedan adaptarse mejor a las personas que no pueden utilizar cómodamente los teclados convencionales. Aunque existen soluciones como el dictado por voz, estas no siempre son adecuadas en todos los contextos, especialmente en espacios compartidos o silenciosos.

Por ello, este proyecto busca aportar una solución desde el ámbito de los sistemas empotrados, combinando el desarrollo técnico con una finalidad práctica y social.


\section{Estructura de la memoria}

Esta sección permite al lector comprender la secuencia lógica de cómo se
ha desarrollado el proyecto o la investigación. Presentará un resumen de
los diferentes capítulos que conforman la memoria (a poder ser en un
único párrafo, como mucho dos).

\chapter{Estado del arte}
\label{ch:estado-del-arte}
1. Qué alternativas existen.
- Teclados ergonómicos.
- Teclados adaptados para personas con movilidad reducida.
- Dictado por voz. 
- Pulsadores o switches adaptados.% NOTE
- Interfaces basadas en gestos y Eye tracking.% NOTE
- Sistemas empotrados usados como dispositivos HID 

2. Cómo funcionan.
3. Qué ventajas tienen.
4. Qué limitaciones tienen.
5. Qué hueco queda para justificar tu proyecto.




\section{Introducción}
\label{sec:estado-arte-introduccion}

En la actualidad, el mercado de dispositivos de entrada de texto es muy amplio. Sin embargo, aunque existen muchas alternativas, el modelo de teclado más extendido alrededor del mundo sigue teniendo ciertas limitaciones.

El teclado convencional moderno obliga al usuario a colocar las manos en una posición concreta y a repartir la carga de escritura de manera desigual entre todos los dedos. Esto puede generar problemas, ya que algunos dedos, como el meñique, tienen menos fuerza y destreza, aun así deben encargarse desproporcionadamente de teclas importantes. Mientras que el dedo más hábil, el pulgar solo se usa para una tecla. Además, el diseño recto del teclado suele hacer que las manos permanezcan juntas y que las muñecas adopten posiciones poco naturales durante periodos prolongados.

Este tipo de uso no tiene por qué afectar a todas las personas de la misma manera, pero sí puede contribuir a situaciones de incomodidad, fatiga o molestias asociadas al uso repetitivo de la mano y la muñeca. De hecho, el síndrome del túnel carpiano está relacionado con la compresión del nervio mediano en la muñeca, y algunos factores como los movimientos repetitivos o las posiciones extremas de la mano y la muñeca pueden aumentar el riesgo de sufrir este tipo de problemas \cite{aaosCarpalTunnel}.

El problema es todavía mayor para las personas que, por las características de su cuerpo o por alguna limitación de movilidad, no pueden utilizar cómodamente un teclado convencional. En estos casos, muchas veces se ven obligadas a recurrir a alternativas que pueden ser incómodas, poco prácticas o ineficientes en determinados contextos. Algunas soluciones no están pensadas realmente para personas que no tienen un cuerpo estándar, por lo que pueden acabar siendo tediosas, difíciles de usar o poco adaptadas a sus necesidades reales.

Con este contexto, en este capítulo se van a explorar distintas alternativas existentes para la entrada de texto y la interacción con el ordenador. El objetivo es estudiar cómo funcionan, qué ventajas ofrecen y qué limitaciones presentan, para poder justificar después la propuesta desarrollada en este trabajo.

\section{Clasificación de alternativas}
\label{sec:clasificacion-alternativas}

% Aquí va la tabla con enlaces
\begin{table}
    \centering
    \caption{Alternativas analizadas en el estado del arte}
    \label{tab:alternativas-estado-arte}
    \begin{tabular}{p{4cm} p{8cm}}
        \toprule
        \textbf{Alternativa} & \textbf{Descripción general} \\
        \midrule
        \hyperref[sec:teclados-ergonomicos]{Teclados ergonómicos} 
        & Modifican la forma física del teclado para mejorar la postura de manos y muñecas. \\

        \hyperref[sec:teclados-adaptados]{Teclados adaptados} 
        & Están diseñados para personas con movilidad reducida o necesidades específicas. \\

        \hyperref[sec:dictado-voz]{Dictado por voz} 
        & Permite introducir texto mediante reconocimiento de voz. \\

        \hyperref[sec:pulsadores-switches]{Pulsadores o switches adaptados} 
        & Permiten generar entradas mediante uno o varios botones físicos adaptados. \\

        \hyperref[sec:interfaces-gestos]{Interfaces basadas en gestos} 
        & Utilizan movimientos del cuerpo, manos o cabeza como método de interacción. \\

        \hyperref[sec:eye-tracking]{Seguimiento ocular} 
        & Permite controlar el sistema mediante la mirada. \\
        \bottomrule
    \end{tabular}
\end{table}


\section{Teclados ergonómicos}
\label{sec:teclados-ergonomicos}

Los teclados ergonómicos son una de las primeras alternativas que aparecen cuando se intenta mejorar la experiencia de escritura respecto a un teclado convencional. No eliminan el teclado como método de entrada, pero modifican su forma, inclinación o distribución para intentar que la postura de las manos, muñecas y brazos sea más natural durante el uso.

\subsection{Cómo funcionan}

Este tipo de teclados suelen basarse en cambiar la geometría del teclado tradicional. Algunos modelos dividen el teclado en dos zonas, una para cada mano, de manera que las muñecas no tengan que permanecer tan juntas ni giradas hacia dentro. Otros incorporan inclinación, curvatura o reposamuñecas para favorecer una posición más cómoda durante periodos prolongados de escritura.

La idea principal es reducir ciertas posturas forzadas que aparecen al escribir durante mucho tiempo en un teclado recto. Por ejemplo, los teclados divididos buscan que cada mano pueda colocarse en una posición más cercana a su postura natural, reduciendo la desviación lateral de la muñeca. También existen diseños con inclinación vertical, cuyo objetivo es disminuir la rotación del antebrazo al escribir \cite{marklin1999splitkeyboards}.

\subsection{Qué ventajas tienen}

La principal ventaja de los teclados ergonómicos es que mantienen una forma de uso parecida a la de un teclado convencional, por lo que el usuario no tiene que aprender un sistema completamente nuevo. Esto los convierte en una alternativa bastante accesible para personas que ya saben escribir con teclado, pero que buscan una postura más cómoda o una reducción de la fatiga.

Además, como se mencionó antes, su principal ventaja es que gracias a su forma más orgánica, no se tienen que forzar tanto las articulaciones para escribir, por lo cual son mucho más cómodos de utilizar que un teclado convencional.
\subsection{Qué limitaciones tienen}

A pesar de sus ventajas, los teclados ergonómicos no resuelven completamente el problema tratado en este trabajo. Aunque mejoran la postura respecto a un teclado convencional, siguen exigiendo que el usuario pueda mover varios dedos, alcanzar muchas teclas distintas y coordinar ambas manos con cierta precisión.

Por tanto, pueden ser útiles para personas que sienten incomodidad con un teclado tradicional, pero no son una solución suficiente para usuarios con movilidad reducida, falta de fuerza en algunos dedos o dificultad para realizar movimientos finos. En estos casos, el problema no es solo la forma del teclado, sino la propia necesidad de utilizar muchas teclas físicas repartidas por una superficie amplia.

Por este motivo, los teclados ergonómicos representan una mejora del teclado convencional, pero no una alternativa radicalmente distinta. Esta limitación es importante para este proyecto, ya que la solución propuesta busca reducir la interacción a un número mucho menor de controles físicos.



\section{Teclados adaptados para personas con movilidad reducida}
\label{sec:teclados-adaptados}

Los teclados adaptados para personas con movilidad reducida son dispositivos pensados para facilitar el acceso al ordenador cuando el uso de un teclado convencional resulta difícil, incómodo o poco preciso. A diferencia de los teclados ergonómicos, que suelen centrarse en mejorar la postura de una persona que ya puede escribir con relativa normalidad, estos dispositivos tienen un enfoque más cercano a la accesibilidad. Es decir, no buscan únicamente mejorar la comodidad, sino permitir que el usuario pueda interactuar con el ordenador de una forma más ajustada a sus capacidades físicas \cite{whoAssistiveTechnology}.

\subsection{Cómo funcionan}

Estos teclados pueden funcionar de formas bastante distintas según el tipo de necesidad que intenten cubrir. Algunos modelos aumentan el tamaño de las teclas para que sea más fácil pulsarlas, mientras que otros reducen el número de teclas disponibles o incorporan protectores que separan unas teclas de otras para evitar pulsaciones accidentales. También existen teclados con alto contraste, distribuciones simplificadas o superficies preparadas para usuarios que no tienen un control fino de los dedos \cite{abilityNetKeyboardMouseAlternatives}.

En estos casos, el objetivo no siempre es escribir más rápido. Muchas veces lo importante es que la escritura sea posible, estable y menos frustrante. Por eso, este tipo de soluciones suelen dar más importancia a la precisión y a la facilidad de pulsación que a la velocidad. Un teclado adaptado puede parecer menos eficiente desde el punto de vista de una persona que usa un teclado convencional sin dificultad, pero para determinados usuarios puede marcar la diferencia entre poder utilizar un ordenador o depender completamente de otra persona.

También es habitual que estos teclados se combinen con otros elementos de apoyo, como reposamanos, soportes, carcasas, protectores de teclado o dispositivos de señalización alternativos. De esta forma, el teclado deja de ser un dispositivo genérico y pasa a formar parte de una solución más personalizada, pensada para la situación concreta del usuario \cite{abilityNetKeyboardMouseAlternatives}.

\subsection{Qué ventajas tienen}

Una de sus ventajas más importantes es que parten de una idea sencilla: no todas las personas pueden interactuar con el ordenador de la misma manera. En lugar de obligar al usuario a adaptarse a un teclado estándar, modifican el propio dispositivo para hacerlo más accesible. Esto encaja con el objetivo general de las tecnologías de asistencia, que buscan mejorar la autonomía, la participación y la calidad de vida de las personas que las necesitan \cite{whoAssistiveTechnology}.

Otra ventaja es que mantienen una lógica parecida a la de un teclado normal. Esto puede facilitar el aprendizaje, porque el usuario no tiene que entender un sistema completamente nuevo. Si una persona ya conoce, aunque sea parcialmente, la distribución de un teclado, un teclado adaptado puede ser una mejora directa frente al teclado convencional.

Además, al tratarse de dispositivos físicos, pueden utilizarse en contextos donde otras alternativas no son tan cómodas. Por ejemplo, no obligan a hablar en voz alta, no dependen de una cámara y no necesitan unas condiciones de iluminación concretas. Esto puede hacerlos útiles en espacios compartidos, aulas, oficinas o bibliotecas.

\subsection{Qué limitaciones tienen}

A pesar de sus ventajas, estos teclados siguen teniendo una limitación importante: continúan basándose en la idea de pulsar teclas físicas distribuidas sobre una superficie. Aunque las teclas sean más grandes, estén mejor separadas o sean más fáciles de distinguir, el usuario todavía necesita cierta movilidad, fuerza y precisión para desplazarse por el teclado.

También pueden ser dispositivos grandes, poco portátiles o demasiado específicos para un perfil concreto de usuario. En accesibilidad, una solución que funciona bien para una persona no tiene por qué servir para otra, porque pequeñas diferencias en movilidad, fuerza o coordinación pueden cambiar completamente la utilidad del dispositivo.

Por último, muchos teclados adaptados mejoran el acceso al teclado convencional, pero no reducen de forma radical la cantidad de acciones físicas necesarias para escribir. Esta limitación es importante en relación con este trabajo, ya que la propuesta desarrollada busca explorar una entrada más reducida, basada en un joystick y dos botones, de forma que la interacción no dependa de alcanzar muchas teclas diferentes.


\section{Dictado por voz}
\label{sec:dictado-voz}

El dictado por voz es una de las alternativas más conocidas al uso del teclado. En lugar de introducir el texto mediante pulsaciones físicas, el usuario habla y el sistema transforma esa señal de audio en palabras escritas. Es una solución interesante porque, en principio, permite escribir sin utilizar las manos, lo que puede ser una ventaja clara para personas que no pueden usar cómodamente un teclado convencional.

\subsection{Cómo funcionan}

Los sistemas de dictado por voz se basan en tecnologías de reconocimiento automático del habla. De forma general, el sistema capta la voz mediante un micrófono, procesa la señal de audio y la convierte en texto. En algunos casos, además de escribir texto, estos sistemas permiten ejecutar comandos, corregir palabras, seleccionar elementos de la pantalla o controlar partes del ordenador mediante instrucciones habladas \cite{webSpeechApi,microsoftVoiceAccess}.

Esta tecnología está presente en muchos sistemas actuales. Por ejemplo, algunos sistemas operativos incorporan funciones de accesibilidad que permiten dictar texto y controlar el equipo con la voz. También existen interfaces web que permiten integrar reconocimiento de voz en aplicaciones, diferenciando entre la parte de reconocimiento del habla y la parte de síntesis de voz \cite{webSpeechApi,appleVoiceControl}.

Desde el punto de vista del usuario, el funcionamiento parece sencillo: se habla y aparece texto en pantalla. Sin embargo, por debajo depende de varios factores, como la calidad del micrófono, el ruido del entorno, la pronunciación, el idioma, el acento o la capacidad del sistema para interpretar correctamente pausas, signos de puntuación y correcciones.

\subsection{Qué ventajas tienen}

La ventaja más evidente del dictado por voz es que reduce o incluso elimina la necesidad de usar las manos para escribir. Esto lo convierte en una alternativa muy útil para personas que tienen dolor, fatiga, lesiones o dificultades de movilidad en manos, muñecas o brazos.

También puede ser una forma de entrada bastante rápida cuando las condiciones son buenas. Para redactar texto largo, hablar puede resultar más natural que escribir, especialmente si el usuario no necesita introducir muchos símbolos, comandos técnicos o correcciones constantes. Además, al estar integrado en muchos dispositivos actuales, no siempre requiere comprar hardware específico.

Otra ventaja importante es que no obliga al usuario a aprender una distribución nueva de teclas. La interacción se basa en el habla, que es una forma de comunicación natural para muchas personas. En ese sentido, el dictado por voz puede ser una solución muy potente cuando el entorno acompaña y cuando la persona puede hablar de forma cómoda y clara.

\subsection{Qué limitaciones tienen}

A pesar de sus ventajas, el dictado por voz no es una solución válida para todos los contextos. Uno de sus problemas principales es que obliga al usuario a hablar en voz alta. Esto puede ser incómodo o directamente inviable en lugares como bibliotecas, aulas, oficinas compartidas o espacios donde hay más personas trabajando alrededor. En esos casos, aunque la tecnología funcione bien, el contexto hace que no sea una alternativa práctica.

También hay situaciones en las que el reconocimiento puede fallar. El ruido de fondo, una mala calidad del micrófono o una pronunciación que el sistema no reconozca bien pueden provocar errores en el texto generado. Además, algunos estudios señalan que los sistemas de reconocimiento de voz no funcionan igual de bien para todas las personas, especialmente en casos de usuarios con diferencias en el habla, como la tartamudez, donde pueden aparecer cortes, malentendidos o predicciones que no representan correctamente lo que la persona quería decir \cite{lea2023stutter}.

Otro aspecto importante es la privacidad. Para dictar texto, el usuario tiene que decir en voz alta lo que quiere escribir. Esto puede ser un problema si se está trabajando con información personal, académica, médica, laboral o simplemente privada. Además, las tecnologías de asistencia basadas en voz pueden plantear preocupaciones de seguridad y privacidad cuando procesan audio del usuario o del entorno \cite{aloufi2021privacyASR}.

Por todo ello, el dictado por voz es una alternativa muy útil en determinados casos, pero no resuelve completamente el problema que se trata en este trabajo. Puede evitar el uso físico del teclado, pero depende mucho del entorno, de la voz del usuario y de que sea aceptable hablar en voz alta. La solución propuesta en este proyecto busca cubrir precisamente otro tipo de situación: permitir la entrada de texto de forma silenciosa, discreta y con una interacción física reducida.



\section{Pulsadores o switches adaptados}
\label{sec:pulsadores-switches}

Los pulsadores o \textit{switches} adaptados son otra alternativa importante dentro de las tecnologías de asistencia. Su idea es bastante sencilla: reducir la interacción a una o varias acciones físicas simples, como pulsar un botón, mover ligeramente una parte del cuerpo o activar un sensor. Esto puede ser útil para personas que no pueden utilizar con precisión un teclado, un ratón o una pantalla táctil.

\subsection{Cómo funcionan}

Un switch adaptado puede entenderse como un dispositivo que envía una señal cuando el usuario realiza una acción concreta. Esa acción puede ser pulsar con la mano, presionar con el pie, mover la cabeza, accionar la barbilla o realizar cualquier otro movimiento que la persona pueda repetir de forma cómoda y fiable \cite{ableNetWhatIsSwitch}.

Normalmente, estos pulsadores no se utilizan solos, sino conectados a un sistema que interpreta la señal. Por ejemplo, pueden conectarse a un comunicador, a un ordenador, a una tableta o a una interfaz de accesibilidad. En muchos casos se usan junto con sistemas de escaneo: el sistema va resaltando elementos de la pantalla uno por uno, o por grupos, y el usuario pulsa el switch cuando aparece la opción que quiere seleccionar \cite{googleSwitchAccess,appleSwitchControlIOS}.

Este funcionamiento permite controlar interfaces complejas utilizando muy pocas entradas físicas. Con un solo pulsador se puede seleccionar una opción cuando el sistema la resalta automáticamente. Con dos pulsadores, uno puede servir para avanzar entre opciones y otro para seleccionar. Este tipo de interacción es más lenta que escribir en un teclado convencional, pero puede abrir el acceso al ordenador a personas para las que otros métodos no son viables.

\subsection{Qué ventajas tienen}

La ventaja principal de los switches adaptados es que no exigen movimientos complejos. En lugar de tener que desplazar las manos por un teclado completo, el usuario solo necesita realizar una acción física concreta y repetible. Esto hace que puedan adaptarse a perfiles muy diferentes, ya que el pulsador puede colocarse donde la persona tenga más control: cerca de la mano, del pie, de la cabeza o de otra parte del cuerpo.

También son dispositivos relativamente simples desde el punto de vista físico. Muchos switches se basan en una entrada binaria, activado o no activado, lo que facilita su integración con sistemas electrónicos y programas de accesibilidad. Esta simplicidad es una ventaja importante, porque permite construir soluciones robustas y fáciles de entender para el usuario.

Otra ventaja es que los sistemas operativos actuales ya incluyen funciones de control mediante switches. Android permite conectar pulsadores externos por USB o Bluetooth y usarlos para seleccionar elementos, desplazarse o introducir texto mediante Switch Access \cite{googleSwitchAccess}. En el caso de Apple, Switch Control permite navegar por elementos de la pantalla mediante escaneo y seleccionar opciones cuando están resaltadas \cite{appleSwitchControlMac}. Esto demuestra que no se trata de una idea aislada, sino de una forma de interacción reconocida dentro de la accesibilidad digital.

\subsection{Qué limitaciones tienen}

La principal limitación de los switches adaptados es la velocidad. Al reducir la interacción a uno o pocos pulsadores, muchas acciones necesitan más pasos que en un teclado o ratón convencional. Si el sistema tiene que ir recorriendo opciones hasta que el usuario selecciona una, escribir una palabra o navegar por una interfaz puede convertirse en un proceso lento.

También requieren una configuración bastante personalizada. No basta con elegir un pulsador cualquiera: hay que decidir dónde colocarlo, qué fuerza necesita, qué movimiento va a usar la persona y qué acción del sistema va a activar. Un switch mal colocado o demasiado duro puede acabar siendo incómodo, cansado o directamente inútil.

Además, aunque los switches reducen mucho la complejidad física de la interacción, a veces trasladan parte de esa complejidad al software. El usuario puede tener que esperar al escaneo, memorizar secuencias o depender de menús intermedios para llegar a la acción que quiere realizar. Por tanto, el problema no desaparece por completo, sino que cambia de forma.

En relación con este trabajo, los switches adaptados son una referencia muy importante, porque muestran que es posible controlar un sistema con pocas acciones físicas. Sin embargo, la propuesta de este TFG intenta ir un paso más allá: mantener esa idea de interacción reducida, pero combinándola con un joystick y dos botones para ofrecer más posibilidades de entrada sin volver a depender de un teclado completo.



\section{Interfaces basadas en gestos}
\label{sec:interfaces-gestos}

Las interfaces basadas en gestos permiten controlar un sistema mediante movimientos del cuerpo, de las manos, de la cabeza o incluso de la cara. En lugar de pulsar teclas físicas, el usuario realiza un gesto y el sistema intenta reconocerlo para convertirlo en una acción. Esta idea resulta interesante en el contexto de la accesibilidad, porque permite plantear formas de interacción que no dependen directamente de un teclado convencional.

\subsection{Cómo funcionan}

Estas interfaces suelen utilizar sensores para detectar el movimiento del usuario. En algunos casos se emplean cámaras convencionales, cámaras de profundidad o sensores de movimiento. En otros, se utilizan dispositivos colocados sobre el cuerpo, como guantes, acelerómetros o sensores inerciales. A partir de esos datos, el sistema interpreta si el usuario ha realizado un gesto concreto y lo asocia a una acción determinada, como mover el cursor, seleccionar un elemento o ejecutar un comando \cite{oudah2020gestureRecognition}.

Un ejemplo de este tipo de interacción son los sistemas que permiten mover el puntero con la cabeza o la cara. En macOS, por ejemplo, la función \textit{Head Pointer} permite controlar el movimiento del cursor usando la cámara del ordenador para detectar el movimiento de la cabeza o de la cara \cite{appleHeadPointerMac}. De forma similar, en iPhone existen funciones de seguimiento de cabeza que permiten mover un puntero en pantalla y realizar acciones manteniendo la posición o usando movimientos faciales concretos \cite{appleHeadTrackingIphone}.

También existen proyectos orientados a accesibilidad que utilizan gestos faciales para controlar el ordenador. Un caso interesante es Project Gameface, desarrollado por Google, que plantea un ratón manos libres basado en movimientos de cabeza y gestos faciales. En este sistema, acciones como levantar las cejas o abrir la boca pueden asignarse a funciones como hacer clic, arrastrar o mover el cursor \cite{googleProjectGameface}.

\subsection{Qué ventajas tienen}

La principal ventaja de las interfaces basadas en gestos es que pueden reducir la dependencia de dispositivos físicos tradicionales. Para una persona que no puede utilizar bien un teclado o un ratón, poder controlar el sistema con movimientos de la cabeza, de la cara o de otra parte del cuerpo puede abrir una forma de interacción más accesible.

Otra ventaja es que los gestos pueden resultar bastante naturales. Mover la cabeza hacia una dirección, levantar las cejas o realizar un movimiento concreto puede ser más intuitivo que aprender combinaciones de teclas o recorrer menús mediante escaneo. Además, al no requerir necesariamente contacto físico con el dispositivo, pueden ser útiles cuando el usuario tiene dificultad para ejercer fuerza, mantener precisión o realizar pulsaciones repetidas.

También es importante que muchas de estas soluciones pueden funcionar con hardware que ya está presente en dispositivos actuales, como una cámara integrada. Esto puede reducir la necesidad de comprar periféricos específicos y facilita que este tipo de interacción pueda integrarse en ordenadores, teléfonos o tabletas.

\subsection{Qué limitaciones tienen}

A pesar de sus posibilidades, las interfaces basadas en gestos tienen limitaciones claras. La primera es que dependen mucho de que el sistema reconozca bien el movimiento. En los sistemas basados en cámara, factores como la iluminación, la posición del usuario, el fondo, la distancia a la cámara o la oclusión de alguna parte del cuerpo pueden afectar al resultado. De hecho, Apple recomienda que la cámara tenga una visión clara de la cara y que esta esté suficientemente iluminada para obtener mejores resultados en el seguimiento de cabeza \cite{appleHeadTrackingIphone}.

Otro problema es que no todos los usuarios pueden realizar los mismos gestos. Una persona puede tener movilidad suficiente para mover la cabeza, pero no para mantener una posición estable; otra puede realizar movimientos faciales, pero cansarse rápido; y otra puede tener movimientos involuntarios que dificulten el reconocimiento. Por eso, aunque los gestos parecen una solución muy flexible, en la práctica necesitan mucha personalización.

También puede aparecer fatiga. Mantener la cabeza en una posición, repetir gestos faciales o realizar movimientos amplios durante mucho tiempo puede resultar cansado. Además, si el sistema interpreta mal un gesto, el usuario puede acabar dedicando más esfuerzo a corregir errores que a realizar la tarea principal.

En relación con este trabajo, las interfaces basadas en gestos muestran que es posible reducir la dependencia del teclado convencional. Sin embargo, también introducen una dependencia fuerte del entorno, de la cámara y de la capacidad del usuario para realizar gestos reconocibles de forma estable. La propuesta desarrollada en este TFG busca una alternativa más física y directa, basada en un joystick y dos botones, que mantenga una interacción reducida sin depender de reconocimiento visual continuo.



\section{Sistemas basados en seguimiento ocular}
\label{sec:eye-tracking}

Los sistemas basados en seguimiento ocular, también conocidos como sistemas de \textit{eye tracking}, permiten controlar una interfaz utilizando la mirada. En este tipo de soluciones, el usuario no necesita pulsar teclas ni mover un ratón de forma directa, sino mirar a una zona de la pantalla para desplazar un puntero, seleccionar un elemento o escribir mediante un teclado virtual.

\subsection{Cómo funcionan}

Estos sistemas suelen utilizar una cámara o un sensor especializado para detectar la posición de los ojos y estimar hacia qué punto de la pantalla está mirando el usuario. Antes de poder utilizarlos, normalmente es necesario realizar una calibración. Durante este proceso, el usuario sigue con la mirada una serie de puntos que aparecen en la pantalla, y el sistema utiliza esa información para relacionar el movimiento ocular con coordenadas concretas de la interfaz \cite{appleEyeTrackingIphone}.

Una vez calibrado, el sistema puede mover un puntero siguiendo la mirada del usuario. Para seleccionar un elemento, muchas soluciones utilizan un mecanismo conocido como \textit{dwell}. Esto consiste en mantener la mirada fija sobre una zona durante un tiempo determinado; cuando ese tiempo se cumple, el sistema interpreta que el usuario quiere hacer clic o activar esa opción \cite{appleEyeTrackingIphone}.

En el ámbito de la accesibilidad, el seguimiento ocular se utiliza especialmente en personas con movilidad muy reducida. Algunos sistemas permiten emular el ratón, seleccionar elementos de Windows, utilizar teclados en pantalla o apoyarse en predicción de palabras para facilitar la escritura \cite{adcetEyeGazeTechnology}. De esta forma, la mirada se convierte en el principal canal de interacción con el ordenador.

\subsection{Qué ventajas tienen}

La ventaja más importante del seguimiento ocular es que puede permitir el uso de un ordenador a personas que no pueden manejar un teclado, un ratón o un pulsador de forma cómoda. En casos de movilidad muy limitada, poder controlar una interfaz únicamente con los ojos puede ser una herramienta muy valiosa para comunicarse, escribir o acceder a contenidos digitales.

Otra ventaja es que no exige fuerza física ni movimientos amplios. El usuario no tiene que desplazar las manos por una superficie ni mantener pulsaciones repetidas. Esto lo diferencia de otras alternativas que, aunque reducen el número de teclas, siguen dependiendo de cierta movilidad en dedos, manos o brazos.

También puede combinarse con teclados virtuales, predicción de texto y sistemas de comunicación aumentativa y alternativa. Esto hace que no sea solo una herramienta de control del cursor, sino también una posible vía de comunicación para personas con grandes limitaciones motoras \cite{borgestig2015eyeGazePerformance}.

\subsection{Qué limitaciones tienen}

A pesar de sus ventajas, el seguimiento ocular también tiene limitaciones importantes. La primera es que depende mucho de la calibración y de las condiciones de uso. Si cambia la posición de la cabeza, la distancia a la pantalla, la iluminación o la colocación del dispositivo, puede ser necesario recalibrar el sistema o ajustar de nuevo la postura del usuario \cite{appleEyeTrackingIphone}.

Otra limitación es la precisión. Mirar a un punto concreto de la pantalla no siempre significa que el usuario quiera seleccionarlo. Este problema se suele conocer como \textit{Midas Touch}: si el sistema interpreta cualquier mirada como una intención de hacer clic, pueden producirse acciones no deseadas. Para evitarlo, muchos sistemas utilizan tiempos de permanencia, pero esto también puede hacer que la interacción sea más lenta \cite{rajanna2018gazeAssisted}.

Además, el uso prolongado de la mirada como método principal de entrada puede producir cansancio. Mantener la vista fija, repetir selecciones con los ojos o escribir en un teclado virtual durante mucho tiempo puede resultar más exigente de lo que parece. Por eso, aunque estos sistemas son muy útiles en algunos casos, no siempre son cómodos para tareas largas o para usuarios que pueden utilizar otras partes del cuerpo con menos esfuerzo.

En relación con este trabajo, el seguimiento ocular representa una alternativa muy avanzada y potente, pero también más compleja. Requiere cámara, calibración, buenas condiciones de uso y una interfaz pensada para reducir errores. La propuesta de este TFG busca una solución más sencilla desde el punto de vista técnico y de uso, basada en un joystick y dos botones, que pueda funcionar de forma silenciosa y sin depender de la mirada ni de una cámara.


\section{Conclusiones del estado del arte}
\label{sec:conclusiones-estado-arte}

A lo largo de este capítulo se han analizado distintas alternativas al teclado convencional. Algunas de ellas, como los teclados ergonómicos, intentan mejorar la postura y reducir la incomodidad durante la escritura, pero siguen manteniendo la misma idea de base: el usuario debe utilizar un conjunto amplio de teclas y coordinar varios dedos para poder escribir.

Otras soluciones, como los teclados adaptados y los pulsadores, se acercan más al ámbito de la accesibilidad. Estas alternativas reconocen que no todas las personas pueden interactuar con el ordenador de la misma manera, y por eso modifican el dispositivo o reducen el número de acciones físicas necesarias. Sin embargo, también presentan limitaciones. En algunos casos siguen necesitando una superficie grande de teclas; en otros, la escritura puede volverse lenta o depender demasiado de sistemas de escaneo.

También existen soluciones que eliminan casi por completo el uso físico del teclado, como el dictado por voz, las interfaces basadas en gestos o el seguimiento ocular. Estas tecnologías pueden ser muy útiles en determinados contextos, pero no siempre son adecuadas para un uso cotidiano. El dictado por voz, por ejemplo, depende del entorno y obliga al usuario a hablar en voz alta. Las interfaces por gestos y los sistemas de seguimiento ocular, por su parte, suelen depender de cámaras, calibración, iluminación y reconocimiento correcto del movimiento.

Por tanto, se puede observar que no existe una única solución válida para todos los usuarios. Cada alternativa resuelve una parte del problema, pero también introduce nuevas limitaciones. Esta idea es importante para el presente trabajo, ya que el objetivo no es sustituir todas las tecnologías existentes, sino explorar una opción diferente dentro del conjunto de soluciones posibles.

A partir de este análisis, la propuesta desarrollada en este TFG se orienta hacia una interacción física reducida, silenciosa y sencilla, basada en un joystick y dos botones. Esta configuración busca disminuir la dependencia de un teclado completo, evitando al mismo tiempo algunas limitaciones de otras alternativas, como la necesidad de hablar en voz alta, depender de una cámara o recorrer grandes superficies de teclas.

En definitiva, el estado del arte muestra que sigue siendo necesario investigar y desarrollar dispositivos de entrada más adaptables. La solución propuesta parte de esa necesidad y plantea un diseño que intenta equilibrar accesibilidad, simplicidad técnica y posibilidad de uso en contextos reales.



\chapter{Análisis y requisitos}
\label{ch:requisitos}

Después de estudiar las distintas formas de introducir texto en un ordenador o dispositivo móvil, se puede observar que todavía existe espacio para métodos alternativos. Las soluciones actuales pueden ser muy útiles en determinados casos, pero suelen tener limitaciones relacionadas con la comodidad, la velocidad, el entorno de uso o las capacidades físicas necesarias.

A partir de este análisis, se plantea aprovechar el pulgar, uno de los dedos con mayor movilidad y precisión de la mano, como elemento principal de la interacción. También se busca reducir al mínimo el número de dedos necesarios, de forma que el dispositivo pueda resultar accesible para un mayor número de personas.

La propuesta se basa en un joystick y dos botones. El joystick permite desplazarse entre las distintas opciones de escritura, mientras que los botones permiten cambiar entre capas y realizar acciones adicionales. Esta idea parte de los teclados virtuales utilizados en consolas, donde el usuario se desplaza carácter a carácter, pero intenta ofrecer una forma de navegación más rápida, cómoda y con posibilidades de ampliación.

El dispositivo deberá poder sostenerse en la palma de una mano y manejarse, idealmente, con el pulgar, el índice y el dedo corazón. También deberá incorporar alimentación mediante batería y una conexión inalámbrica que permita utilizarlo con distintos tipos de dispositivos.

\section{Descripción del problema}
\label{sec:descripcion-problema}

El problema principal que se intenta abordar en este trabajo es que, aunque existen muchas formas diferentes de introducir texto en un dispositivo, ninguna reúne todas las ventajas deseables. Algunas soluciones son rápidas, pero requieren una movilidad elevada; otras son accesibles, pero pueden resultar lentas, incómodas o poco prácticas en determinados contextos.

Esto provoca que muchos dispositivos estén dirigidos a perfiles de usuario muy concretos. Una solución puede funcionar bien para una persona y no ser adecuada para otra con una movilidad, fuerza o coordinación diferentes.

En este trabajo se plantea la posibilidad de reducir esa separación mediante un dispositivo que sea cómodo, portátil y eficiente, sin estar dirigido únicamente a una demografía concreta. La intención es que pueda ser utilizado tanto por personas que manejan sin dificultad otros dispositivos como por usuarios que carecen de la movilidad necesaria para utilizar cómodamente un teclado tradicional.

La propuesta no pretende ser una solución universal, ya que las capacidades de cada persona son diferentes. Su objetivo es explorar una alternativa que reduzca el número de movimientos y controles físicos necesarios para escribir.

\section{Perfil de usuario}
\label{sec:perfil-usuario}

El objetivo es desarrollar un teclado que pueda ser utilizado por un espectro amplio de la sociedad, prestando especial atención a grupos que suelen estar poco representados en el diseño de dispositivos de entrada.

El concepto está pensado para poder sujetarse y utilizarse con una sola mano. En su configuración principal, el usuario deberá disponer de movilidad suficiente para manejar el joystick con el pulgar y accionar dos botones con otros dos dedos, idealmente el índice y el corazón.

También será necesario un periodo de aprendizaje. Al tratarse de una forma diferente de escribir, el usuario tendrá que acostumbrarse a la distribución de los caracteres y a las acciones asociadas al joystick y a los botones. Esta curva de aprendizaje es comparable, en cierta medida, a la que existe al aprender mecanografía: al principio la escritura puede ser lenta, pero la práctica debería permitir realizar las acciones de forma más natural y eficiente.

Aunque el primer prototipo se diseñará para utilizarse con tres dedos, la idea de combinar un joystick y dos pulsadores podría adaptarse en el futuro a otras posiciones o partes del cuerpo.

\section{Escenario de uso}
\label{sec:escenario-uso}

Uno de los principales beneficios buscados es que el dispositivo pueda utilizarse en situaciones en las que otras alternativas presentan dificultades. Por ejemplo, los sistemas de dictado por voz requieren que el usuario pueda hablar en voz alta y que el ruido del entorno no afecte al reconocimiento. Esto limita su uso en bibliotecas, aulas, oficinas compartidas o lugares en los que se necesita mantener silencio.

Los teclados convencionales también ocupan una superficie considerable y suelen necesitar una mesa o un apoyo estable. Además, obligan al usuario a mantener ambas manos sobre el teclado durante la escritura.

La propuesta de este trabajo busca evitar estas limitaciones mediante un dispositivo pequeño que pueda sostenerse en una sola mano. Esto permitiría utilizarlo sentado, de pie o en situaciones en las que no se dispone de una superficie de apoyo.

Incluso podría utilizarse con la mano dentro de un bolsillo, por ejemplo durante un día frío, siempre que el diseño permita reconocer los controles mediante el tacto. En un aula o una reunión, el usuario podría mantener una postura más libre y escribir sin necesitar una mesa completa.

La idea principal es desarrollar un teclado portátil, silencioso y fácil de transportar, que pueda utilizarse con comodidad en distintos contextos.

\section{Requisitos funcionales}
\label{sec:requisitos-funcionales}

Los requisitos funcionales describen las acciones que debe ser capaz de realizar el dispositivo:

\begin{itemize}
    \item \textbf{RF-01}: El dispositivo debe poder conectarse de forma inalámbrica a un ordenador, teléfono móvil o tableta.

    \item \textbf{RF-02}: El dispositivo debe ser reconocido por el sistema como un teclado.

    \item \textbf{RF-03}: El joystick debe detectar ocho direcciones diferentes.

    \item \textbf{RF-04}: Cada dirección del joystick debe representar un carácter o una acción diferente según la capa activa.

    \item \textbf{RF-05}: Los dos botones deben permitir cambiar entre las distintas capas de caracteres.

    \item \textbf{RF-06}: El teclado principal debe permitir introducir todas las letras del alfabeto.

    \item \textbf{RF-07}: El sistema debe permitir alternar entre letras minúsculas y mayúsculas.

    \item \textbf{RF-08}: Una de las opciones disponibles debe permitir abrir un teclado virtual secundario.

    \item \textbf{RF-09}: El teclado virtual secundario debe permitir introducir números, signos de puntuación y caracteres especiales.

    \item \textbf{RF-10}: El dispositivo debe ofrecer información visual sobre la capa activa y la selección realizada.

    \item \textbf{RF-11}: El dispositivo debe permitir que con una combinación de botones se entre un modo de lectura para ver los caracteres pero que no escriba nada.
\end{itemize}

\section{Requisitos no funcionales}
\label{sec:requisitos-no-funcionales}

Los requisitos no funcionales definen las condiciones de calidad y uso que debe cumplir el prototipo:

\begin{itemize}
    \item \textbf{RNF-01}: El dispositivo debe tener un tamaño suficientemente reducido para poder sostenerse en la palma de una mano.

    \item \textbf{RNF-02}: Debe poder manejarse con una sola mano y, en su configuración principal, con tres dedos.

    \item \textbf{RNF-03}: La interacción debe ser silenciosa para permitir su uso en espacios compartidos.

    \item \textbf{RNF-04}: El dispositivo debe ser ligero y fácil de transportar.

    \item \textbf{RNF-05}: Los controles deben poder distinguirse con facilidad mediante el tacto.

    \item \textbf{RNF-06}: El coste de los componentes debe mantenerse reducido.

    \item \textbf{RNF-07}: El sistema debe responder a las acciones del usuario sin retardos que dificulten la escritura.

    \item \textbf{RNF-08}: El funcionamiento básico debe poder aprenderse sin necesidad de conocimientos técnicos.
\end{itemize}

\section{Restricciones del proyecto}
\label{sec:restricciones-proyecto}

Una de las principales restricciones del proyecto es el tamaño. El objetivo de introducir todo el dispositivo en la palma de una mano limita el espacio disponible para colocar la batería, el microcontrolador, el joystick, los botones, la matriz de LED y el resto de las conexiones.

Esta limitación afecta especialmente a la distribución interna de los componentes, al tamaño de la batería y a la autonomía que puede alcanzar el prototipo. También obliga a diseñar una carcasa que aproveche bien el espacio sin dejar de ser cómoda de sujetar.

Para el desarrollo se utilizará un microcontrolador ESP32 debido a su tamaño, capacidad de procesamiento y conectividad Bluetooth. El dispositivo contará con un joystick, dos botones y una matriz de LED para mostrar información sobre la entrada seleccionada. La carcasa será fabricada mediante impresión 3D.

La experiencia limitada en diseño e impresión 3D también supone una restricción, ya que será necesario realizar varias versiones de la carcasa hasta encontrar una forma que permita colocar correctamente los componentes y que resulte cómoda durante el uso.

Por último, el proyecto se centra en desarrollar y validar un prototipo funcional. Por ello, aspectos como la fabricación industrial, la certificación del producto o la realización de pruebas clínicas con usuarios quedan fuera del alcance de esta primera versión.



\chapter{Diseño de la solución}
\label{ch:diseno}

En este capítulo se describe el diseño planteado para el dispositivo antes de entrar en los detalles concretos de su construcción. El objetivo principal ha sido reducir un teclado completo a una interfaz formada por un joystick y dos botones, manteniendo la posibilidad de introducir letras, números, símbolos y algunas acciones habituales de escritura.

La solución se ha diseñado para poder utilizarse con una sola mano. El pulgar se encarga de mover el joystick, mientras que los botones pueden accionarse con los dedos índice y corazón. Para ampliar la cantidad de caracteres disponibles sin añadir más controles físicos, el sistema se organiza mediante varias capas.

\section{Arquitectura general del sistema}
\label{sec:arquitectura-general}

El sistema está formado por cinco bloques principales: los controles de entrada, el microcontrolador, el sistema de visualización, la comunicación inalámbrica y la alimentación.

El joystick y los dos botones recogen las acciones realizadas por el usuario. Estas señales son procesadas por un microcontrolador XIAO ESP32C3, encargado de determinar la dirección seleccionada, identificar la capa activa y obtener el carácter correspondiente. El carácter se muestra en la matriz de LED y, cuando el dispositivo se encuentra en modo de escritura, se envía al equipo conectado mediante Bluetooth.

La arquitectura general puede resumirse mediante el siguiente flujo:

\begin{center}
    Joystick y botones
    $\longrightarrow$
    XIAO ESP32C3
    $\longrightarrow$
    Bluetooth HID
    $\longrightarrow$
    ordenador, tableta o teléfono
\end{center}

La matriz RGB funciona como salida visual del sistema y permite mostrar el carácter seleccionado, la capa activa y distintos estados del dispositivo.

\begin{table}[ht]
    \centering
    \caption{Bloques principales del sistema}
    \label{tab:bloques-sistema}
    \begin{tabularx}{\textwidth}{@{}lX@{}}
        \toprule
        \textbf{Bloque} & \textbf{Función} \\
        \midrule
        Joystick & Seleccionar una de las ocho direcciones disponibles. \\
        Botones & Elegir la capa activa mediante sus distintas combinaciones. \\
        XIAO ESP32C3 & Procesar las entradas, gestionar las capas y controlar el resto del sistema. \\
        Matriz RGB & Mostrar el carácter seleccionado y ofrecer información visual al usuario. \\
        Bluetooth HID & Enviar las pulsaciones como si procedieran de un teclado convencional. \\
        Alimentación & Proporcionar energía al microcontrolador, los controles y la matriz. \\
        \bottomrule
    \end{tabularx}
\end{table}

\section{Diseño del sistema de escritura}
\label{sec:sistema-escritura}

La parte principal del proyecto es el sistema utilizado para convertir un número reducido de controles físicos en un conjunto completo de caracteres. Para conseguirlo, se combinan las ocho direcciones del joystick con ocho capas diferentes.

Cada dirección representa un carácter o una acción. Los dos botones determinan qué capa se encuentra activa, mientras que una dirección reservada permite alternar entre el banco principal de letras y el banco secundario de números y símbolos.

\subsection{Direcciones del joystick}
\label{subsec:direcciones-joystick}

El joystick proporciona dos señales analógicas, una para el eje horizontal y otra para el eje vertical. Estas señales son leídas mediante el convertidor analógico-digital, o ADC, del microcontrolador. El ADC no identifica directamente una dirección, sino que convierte el nivel de tensión de cada eje en un valor numérico que después debe ser interpretado por el programa \cite{espressifADCOneshot}.

Combinando los valores de ambos ejes se distinguen nueve estados: una posición central y ocho direcciones alrededor de ella. Las direcciones utilizadas son arriba, arriba-derecha, derecha, abajo-derecha, abajo, abajo-izquierda, izquierda y arriba-izquierda.

Alrededor de la posición central se define una zona muerta. Mientras los valores de ambos ejes permanezcan dentro de esta zona, el sistema considera que el joystick no está siendo accionado. Esta medida evita que pequeñas variaciones eléctricas o mecánicas provoquen la escritura accidental de caracteres.

Cuando el joystick sale de la zona central, las lecturas de los dos ejes se comparan con los rangos establecidos para determinar la dirección seleccionada. La selección se considera finalizada cuando el joystick vuelve al centro, lo que permite diferenciar una acción de la siguiente.

\subsection{Organización por capas}
\label{subsec:organizacion-capas}

El teclado está organizado en ocho capas. Las cuatro primeras forman el banco principal y contienen las letras del alfabeto y las acciones más habituales. Las cuatro restantes forman un banco secundario destinado a números, signos de puntuación y caracteres especiales.

Cada capa contiene ocho posiciones, una por cada dirección del joystick. La distribución propuesta para el banco principal se muestra en la Tabla~\ref{tab:capas-principales}.

\begin{table}[ht]
    \centering
    \caption{Distribución de las capas principales}
    \label{tab:capas-principales}
    \resizebox{\textwidth}{!}{%
        \begin{tabular}{lcccccccc}
            \toprule
            \textbf{Capa} &
            \textbf{Arriba} &
            \textbf{Sup. dcha.} &
            \textbf{Derecha} &
            \textbf{Inf. dcha.} &
            \textbf{Abajo} &
            \textbf{Inf. izq.} &
            \textbf{Izquierda} &
            \textbf{Sup. izq.} \\
            \midrule
            0 & e & a & o & s & r & n & i & d \\
            1 & l & c & u & m & p & t & b & g \\
            2 & v & y & q & h & f & z & j & ñ \\
            3 & x & k & w & Intro & Borrar & Mayúsculas & Espacio & Banco secundario \\
            \bottomrule
        \end{tabular}%
    }
\end{table}

Las letras más habituales se han colocado en las primeras capas, de modo que puedan alcanzarse sin mantener pulsados ambos botones al mismo tiempo. La cuarta capa contiene las letras restantes y varias acciones que no necesitan repetirse con tanta frecuencia como una letra normal.

La distribución puede modificarse desde el software sin cambiar el hardware. Esto deja abierta la posibilidad de probar otras organizaciones en el futuro y comparar cuál resulta más rápida o fácil de aprender.

\subsection{Función de los botones}
\label{subsec:funcion-botones}

Los dos botones funcionan como selectores de capa. Cada botón puede encontrarse pulsado o liberado, por lo que la combinación de ambos genera cuatro estados distintos.

\begin{table}[ht]
    \centering
    \caption{Selección de capa mediante los dos botones}
    \label{tab:seleccion-capas}
    \begin{tabular}{ccc}
        \toprule
        \textbf{Botón 1} & \textbf{Botón 2} & \textbf{Capa seleccionada} \\
        \midrule
        Liberado & Liberado & Capa 0 \\
        Liberado & Pulsado  & Capa 1 \\
        Pulsado  & Liberado & Capa 2 \\
        Pulsado  & Pulsado  & Capa 3 \\
        \bottomrule
    \end{tabular}
\end{table}

Esta forma de selección permite acceder a cuatro capas sin incorporar cuatro botones diferentes. Además, la capa activa depende directamente del estado físico de los botones, por lo que el usuario puede volver a la capa inicial simplemente liberándolos.

El joystick HW-504 también incorpora un pulsador al presionarlo hacia dentro. En el diseño actual, la combinación de este pulsador con los otros dos botones se reserva para alternar entre el modo de escritura y el modo de exploración. En el modo de exploración se puede consultar la letra asociada a una dirección sin enviarla al dispositivo conectado.

\subsection{Banco secundario de números y caracteres especiales}
\label{subsec:teclado-secundario}

Una de las posiciones de la cuarta capa está reservada para alternar entre el banco principal y el banco secundario. Esta acción no abre un teclado del sistema operativo, sino que cambia internamente el conjunto de capas utilizado por el dispositivo.

Cuando el banco secundario está activo, las mismas combinaciones de botones seleccionan otras cuatro capas. De esta forma, el funcionamiento físico no cambia: el usuario continúa utilizando los dos botones para elegir una capa y el joystick para seleccionar una de sus ocho posiciones.

\begin{table}[ht]
    \centering
    \caption{Distribución del banco secundario}
    \label{tab:capas-secundarias}
    \resizebox{\textwidth}{!}{%
        \begin{tabular}{lcccccccc}
            \toprule
            \textbf{Capa} &
            \textbf{Arriba} &
            \textbf{Sup. dcha.} &
            \textbf{Derecha} &
            \textbf{Inf. dcha.} &
            \textbf{Abajo} &
            \textbf{Inf. izq.} &
            \textbf{Izquierda} &
            \textbf{Sup. izq.} \\
            \midrule
            4 & 0 & 1 & 2 & 3 & 4 & 5 & 6 & 7 \\
            5 & 8 & 9 & . & : & , & ; & - & \_ \\
            6 & ¿ & ? & ¡ & ! & `` & @ & ( & ) \\
            7 & = & + & * & \& & \textless & \textgreater & Espacio & Banco principal \\
            \bottomrule
        \end{tabular}%
    }
\end{table}

El acceso al banco secundario se ha colocado en la misma posición tanto para entrar como para salir. Esto evita que el usuario tenga que aprender dos acciones diferentes para cambiar entre letras y símbolos.

\section{Diseño hardware}
\label{sec:diseno-hardware}

\subsection{Microcontrolador}
\label{subsec:microcontrolador}

El microcontrolador elegido es el Seeed Studio XIAO ESP32C3. Esta placa utiliza un ESP32-C3 con procesador RISC-V, conectividad Bluetooth Low Energy y varias entradas digitales y analógicas. Sus dimensiones son de aproximadamente \(21 \times 17{,}8\) mm, lo que permite integrarla en una carcasa pensada para sostenerse con una sola mano \cite{xiaoESP32C3}.

La disponibilidad de entradas ADC permite leer los dos ejes analógicos del joystick, mientras que los GPIO restantes se utilizan para los botones, el pulsador del joystick y la matriz RGB. La conectividad Bluetooth integrada evita tener que añadir un módulo inalámbrico separado.

Otro motivo para elegir esta placa es su compatibilidad con ESP-IDF, que permite trabajar directamente con los controladores del ADC, los GPIO, FreeRTOS y la pila Bluetooth. Esto facilita dividir el programa en módulos y controlar con mayor precisión el comportamiento del dispositivo.

\subsection{Joystick y botones}
\label{subsec:joystick-botones}

El joystick HW-504 está formado por dos potenciómetros colocados de forma perpendicular. Uno mide el desplazamiento horizontal y el otro el vertical. Sus salidas se conectan a dos entradas analógicas del microcontrolador.

Además de los dos ejes, el joystick incorpora un pulsador que se activa al presionar la palanca. Este pulsador se utiliza para funciones menos frecuentes, de forma que no interfiera con la escritura normal.

Los dos botones externos se conectan a entradas digitales. Se ha planteado utilizar las resistencias internas de \textit{pull-up} del microcontrolador, por lo que las entradas permanecen en nivel alto cuando los botones están liberados y pasan a nivel bajo al pulsarlos.

La colocación física de los botones es importante. Deben poder accionarse de forma independiente, pero también simultáneamente, ya que la cuarta capa requiere mantener pulsados ambos. También deben ofrecer suficiente respuesta táctil para que el usuario reconozca la pulsación sin necesidad de mirar el dispositivo.

\subsection{Sistema de visualización}
\label{subsec:sistema-visualizacion}

Para mostrar la selección se utiliza una matriz RGB de \(6 \times 10\) diseñada para la familia XIAO. La matriz contiene 60 LED WS2812 controlables de forma individual y recibe los datos mediante una única línea conectada al pin D0 \cite{seeedRGBMatrix}.

La matriz permite representar caracteres mediante una fuente de \(6 \times 10\) píxeles, aprovechando toda su superficie. Al mover el joystick se muestra la letra o símbolo correspondiente a la dirección actual.

También se utilizan distintos colores para ayudar a identificar las capas. De esta forma, el usuario puede relacionar visualmente cada color con una combinación de botones. El brillo debe mantenerse limitado, tanto para evitar molestias al mirar la matriz como para reducir el consumo eléctrico.

La pantalla no pretende mostrar frases completas. Su función es ofrecer una confirmación inmediata de la selección y ayudar durante el proceso de aprendizaje del teclado.

\subsection{Alimentación}
\label{subsec:alimentacion}

El dispositivo está pensado para funcionar mediante una batería recargable de ion de litio o polímero de litio de una celda. El uso de una batería permite mantener el carácter portátil del prototipo y evita depender de una conexión por cable durante la escritura.

La matriz RGB necesita una alimentación estable próxima a 5 V. Por este motivo, se plantea utilizar el módulo PowerBoost 1000C de Adafruit, capaz de elevar la tensión de una batería de 3,7 V hasta aproximadamente 5,2 V y de cargar la batería mientras el sistema continúa alimentado \cite{adafruitPowerBoost1000C}.

La salida del módulo debe alimentar tanto la matriz como el microcontrolador, compartiendo una referencia de tierra común. La conexión a la entrada de 5 V del XIAO deberá realizarse respetando las protecciones recomendadas por el fabricante.

La capacidad definitiva de la batería dependerá del espacio interior disponible y del consumo obtenido durante las pruebas. La matriz será previsiblemente el elemento con mayor influencia sobre la autonomía, por lo que el brillo y el número de LED encendidos deberán mantenerse bajo control.

\subsection{Comunicación Bluetooth}
\label{subsec:comunicacion-bluetooth}

La comunicación se realiza mediante Bluetooth Low Energy utilizando la pila NimBLE disponible en ESP-IDF. El dispositivo actúa como un periférico HID, es decir, como un dispositivo de interfaz humana reconocido por el equipo receptor como un teclado.

NimBLE permite implementar un servidor GATT con el servicio HID, las características necesarias para los informes de teclado y los mecanismos de emparejamiento. ESP-IDF incluye soporte configurable para servicios HID, información del dispositivo y nivel de batería dentro de esta pila \cite{espressifNimBLE}.

Al anunciarse mediante Bluetooth, el prototipo utiliza el nombre \enquote{Teclado Unimano}. Una vez emparejado, el sistema receptor puede recibir pulsaciones sin necesitar una aplicación específica, siempre que sea compatible con teclados Bluetooth HID.

\section{Diseño software}
\label{sec:diseno-software}

El software se ha dividido en varios bloques para separar la lectura de controles, la gestión del sistema de escritura, la visualización y la comunicación Bluetooth.

Esta división evita concentrar todo el comportamiento en un único bucle y facilita modificar cada parte sin afectar directamente al resto.

\subsection{Lectura de entradas}
\label{subsec:lectura-entradas}

Los dos ejes del joystick se leen mediante el controlador ADC en modo \textit{oneshot}. Este modo realiza una conversión bajo demanda cada vez que el programa necesita conocer la posición del joystick \cite{espressifADCOneshot}.

A partir de las dos lecturas se determina si el joystick se encuentra centrado o dentro de una de las ocho direcciones. Para reducir errores se aplican márgenes alrededor de los valores esperados y se exige que una dirección permanezca estable durante varias lecturas antes de aceptarla.

Los botones se gestionan mediante interrupciones GPIO. Cuando se produce una pulsación o liberación, la interrupción envía un evento a una cola de FreeRTOS. Una tarea separada recibe el evento, espera un breve tiempo para eliminar rebotes y lee el estado conjunto de los controles.

El uso de una cola permite mantener la rutina de interrupción lo más corta posible y trasladar el procesamiento a una tarea normal.

\subsection{Gestión de capas}
\label{subsec:gestion-capas}

La capa principal se representa mediante un valor entre 0 y 3. Los estados de los dos botones se interpretan como dos bits: uno de ellos representa el bit de mayor peso y el otro el de menor peso.

Por otro lado, una variable de desplazamiento permite seleccionar el banco de capas. Cuando el desplazamiento vale 0 se utilizan las capas de letras; cuando vale 4 se utilizan las capas de números y símbolos.

La capa real se obtiene mediante:

\begin{equation}
    C_{\text{real}} = C_{\text{botones}} + D_{\text{banco}}
\end{equation}

donde \(C_{\text{botones}}\) puede tomar valores entre 0 y 3, y \(D_{\text{banco}}\) puede valer 0 o 4. El resultado siempre se encuentra entre 0 y 7.

\subsection{Selección y envío de caracteres}
\label{subsec:seleccion-envio}

Los caracteres están almacenados en una tabla de ocho filas y ocho columnas. Las filas representan las capas y las columnas representan las direcciones del joystick.

Cuando se detecta una dirección, el programa calcula la capa real y consulta la posición correspondiente de la tabla. El resultado puede ser una letra, un número, un símbolo o una acción especial, como espacio, borrar, Intro, bloqueo de mayúsculas o cambio de banco.

Antes de enviar el carácter se muestra su representación en la matriz RGB. Si el dispositivo está en modo de escritura, el código HID se envía mediante Bluetooth. Si está en modo de exploración, únicamente se muestra el carácter.

También se ha previsto la repetición de determinadas teclas al mantener el joystick en una dirección durante un tiempo. Esta función se reserva principalmente para espacio, borrado e Intro, ya que repetir automáticamente cualquier letra podría producir errores durante el aprendizaje.

\subsection{Funcionamiento como teclado Bluetooth}
\label{subsec:funcionamiento-bluetooth}

El dispositivo define un descriptor HID de teclado. Cada pulsación se transmite mediante un informe de ocho bytes formado por un byte de modificadores, un byte reservado y seis posiciones para códigos de tecla.

Para simular una pulsación completa se envía primero un informe con la tecla activa y, después de un intervalo breve, otro informe vacío. El segundo informe representa la liberación de la tecla y evita que el sistema receptor interprete que permanece pulsada.

Los caracteres que necesitan modificadores, como las mayúsculas o algunos símbolos, se representan combinando el código de la tecla con modificadores como \textit{Shift} o \textit{AltGr}.

El sistema también gestiona el emparejamiento y conserva la información necesaria en la memoria no volátil del microcontrolador. Si se pierde la conexión, el dispositivo vuelve a anunciarse para permitir una nueva conexión.

\section{Diseño físico y ergonómico}
\label{sec:diseno-fisico}

\subsection{Distribución de los controles}
\label{subsec:distribucion-controles}

La forma general está inspirada en mandos de una sola mano, como el Nunchuk de Wii. Este tipo de distribución permite rodear la carcasa con la mano y mantener el pulgar sobre un control direccional.

El joystick se coloca en la parte superior, dentro del alcance natural del pulgar. Los dos botones se sitúan de forma que puedan accionarse con los dedos índice y corazón sin cambiar el agarre.

La matriz debe permanecer visible mientras se utiliza el dispositivo, pero no debe obligar al usuario a girar excesivamente la muñeca. Su posición también debe evitar que quede tapada por la mano.

\subsection{Dimensiones del dispositivo}
\label{subsec:dimensiones-dispositivo}

El objetivo es que la carcasa pueda sostenerse dentro de la palma de una mano. Sin embargo, las dimensiones finales no dependen únicamente de la comodidad, sino también del espacio ocupado por el joystick, la matriz, el microcontrolador, el sistema de alimentación y la batería.

Por este motivo, las medidas definitivas se determinarán después de colocar todos los componentes en el modelo tridimensional y fabricar los primeros prototipos.

% TODO: añadir las dimensiones finales cuando la carcasa esté terminada.
% Dimensiones aproximadas: ancho x alto x profundidad.

\subsection{Diseño de la carcasa}
\label{subsec:diseno-carcasa}

La carcasa se diseña mediante modelado 3D y se fabrica utilizando impresión 3D. Esto permite modificar rápidamente la posición de los controles y repetir el prototipo hasta encontrar una forma cómoda.

Se plantea dividir la carcasa en dos piezas. La parte frontal contiene el joystick y la matriz, mientras que la parte trasera permite cerrar el dispositivo y sujetar la batería y el resto de los componentes.

La unión entre ambas partes puede realizarse mediante pestañas o encajes, evitando el uso de tornillos siempre que la resistencia obtenida sea suficiente. El interior debe incluir soportes que impidan que las placas y la batería se desplacen durante el uso.

También debe mantenerse accesible el conector de carga, el interruptor de encendido y, durante el desarrollo, el puerto utilizado para programar el microcontrolador.

\section{Decisiones de diseño}
\label{sec:decisiones-diseno}

La decisión principal ha sido utilizar un joystick y dos botones en lugar de un conjunto amplio de teclas. Esto reduce el número de controles físicos, aunque introduce la necesidad de aprender una distribución por capas.

El joystick permite obtener ocho selecciones con un único dedo. Los botones multiplican estas posibilidades sin aumentar demasiado el tamaño del dispositivo. La combinación de ocho direcciones y ocho capas proporciona hasta 64 posiciones, suficientes para incluir letras, números, símbolos y acciones especiales.

La matriz RGB se ha incorporado porque el sistema de capas no resulta evidente durante los primeros usos. Mostrar el carácter seleccionado ayuda al aprendizaje y permite comprobar la entrada antes de enviarla.

La comunicación BLE HID se ha elegido para evitar una aplicación intermedia y permitir que el dispositivo se comporte como un teclado reconocible por distintos sistemas. El XIAO ESP32C3 reúne en una placa de pequeño tamaño las entradas analógicas, los GPIO y la conectividad necesarios para implementar esta solución.

Por último, la fabricación mediante impresión 3D permite adaptar la carcasa a la forma de la mano y modificarla durante el desarrollo. Esta flexibilidad es especialmente importante en un proyecto donde la comodidad y la posición de los controles forman parte del funcionamiento del dispositivo.


\chapter{Implementación}
\label{ch:implementacion}

En este capítulo se describe el proceso seguido para convertir el diseño planteado anteriormente en un prototipo funcional. La implementación se ha dividido en dos partes: por un lado, el montaje y la conexión de los componentes físicos; por otro, el desarrollo del firmware encargado de interpretar las acciones del usuario, mostrar información en la matriz y enviar las pulsaciones mediante Bluetooth.

\section{Montaje del hardware}
\label{sec:montaje-hardware}

El prototipo está formado por un microcontrolador XIAO ESP32C3, un joystick HW-504, dos botones externos y una matriz RGB de \(6 \times 10\) LED. Durante la fase de desarrollo, los componentes se conectaron de forma provisional para poder comprobar su funcionamiento antes de realizar la soldadura y el montaje definitivo dentro de la carcasa.

Las conexiones utilizadas en el firmware se muestran en la Tabla~\ref{tab:conexiones-hardware}.

\begin{table}[ht]
    \centering
    \caption{Conexiones entre los componentes y el XIAO ESP32C3}
    \label{tab:conexiones-hardware}
    \begin{tabularx}{\textwidth}{@{}lXl@{}}
        \toprule
        \textbf{Componente} & \textbf{Señal} & \textbf{GPIO} \\
        \midrule
        Joystick & Eje horizontal & GPIO 2 \\
        Joystick & Eje vertical & GPIO 3 \\
        Joystick & Pulsador integrado & GPIO 4 \\
        Botón 1 & Selección de capa & GPIO 5 \\
        Botón 2 & Selección de capa & GPIO 6 \\
        Matriz RGB & Línea de datos & GPIO 10 \\
        \bottomrule
    \end{tabularx}
\end{table}

Los ejes del joystick se conectan a entradas compatibles con el convertidor analógico-digital del ESP32C3. Los dos botones externos y el pulsador del joystick se conectan a entradas digitales configuradas con resistencias internas de \textit{pull-up}. Por este motivo, una entrada presenta nivel alto mientras el botón está liberado y nivel bajo cuando se pulsa.

Todos los componentes deben compartir una referencia de tierra común. La matriz RGB requiere una alimentación con mayor capacidad de corriente que el resto del circuito, por lo que la distribución de alimentación deberá realizarse evitando que toda la corriente de los LED atraviese el microcontrolador.

% TODO: Completar cuando se realice el montaje definitivo:
% - Tipo y capacidad de la batería utilizada.
% - Conexión final del PowerBoost 1000C.
% - Interruptor de encendido.
% - Fotografías del montaje antes y después de soldar.

\section{Desarrollo del firmware}
\label{sec:desarrollo-firmware}

El firmware se ha desarrollado en lenguaje C utilizando ESP-IDF. Se ha elegido este entorno porque permite acceder directamente a los controladores del ESP32C3, como el ADC, los GPIO, la pila Bluetooth NimBLE y el sistema operativo FreeRTOS.

El programa se ha dividido en varios módulos. El archivo principal contiene la lectura de entradas, la selección de caracteres y la gestión de capas. La comunicación Bluetooth se encuentra separada en un módulo dedicado, mientras que la matriz RGB y la fuente gráfica se gestionan desde otros archivos.

Durante el arranque, la función principal realiza las siguientes operaciones:

\begin{enumerate}
    \item Inicializa la memoria no volátil necesaria para guardar la información de emparejamiento Bluetooth.
    \item Configura las dos entradas ADC utilizadas por el joystick.
    \item Configura los botones y el pulsador del joystick como entradas digitales.
    \item Inicializa la comunicación Bluetooth HID.
    \item Inicializa la matriz RGB.
    \item Crea la cola utilizada para comunicar las interrupciones de los botones.
    \item Crea las tareas de FreeRTOS encargadas de gestionar el joystick y los botones.
    \item Activa las interrupciones asociadas a los controles digitales.
\end{enumerate}

\subsection{Lectura del joystick}
\label{subsec:implementacion-joystick}

El joystick HW-504 genera dos señales analógicas, una para cada eje. Estas señales se leen utilizando el controlador ADC en modo \textit{oneshot}, que permite solicitar una conversión cada vez que el programa necesita conocer la posición.

En las pruebas realizadas se obtuvieron valores centrales aproximados de 1995 para un eje y 1882 para el otro. Los extremos se encuentran próximos a 0 y 4095, ya que el ADC está configurado con una resolución de 12 bits.

A partir de estos valores, la función de lectura clasifica la posición en uno de nueve estados: las ocho direcciones posibles o la posición central. Para cada dirección se define un margen de tolerancia de 300 unidades alrededor de los valores esperados.

La posición central se trata como un estado independiente. Cuando el joystick vuelve al centro, el programa considera que la selección anterior ha finalizado y queda preparado para detectar un nuevo carácter.

Este mecanismo también evita que pequeñas variaciones en las lecturas analógicas produzcan pulsaciones accidentales. Además, antes de aceptar una dirección nueva se exige que permanezca estable durante varias lecturas consecutivas.

La tarea encargada del joystick se ejecuta aproximadamente cada 20 milisegundos. Esta frecuencia permite detectar los movimientos con rapidez sin mantener ocupado constantemente el procesador.

\subsection{Gestión de los botones y las capas}
\label{subsec:implementacion-capas}

Los dos botones externos se utilizan para seleccionar una de las cuatro capas disponibles dentro de cada banco. El estado de ambos se interpreta como un número binario de dos bits:

\begin{equation}
    C = 2B_1 + B_2
\end{equation}

donde \(B_1\) y \(B_2\) representan el estado de cada botón y \(C\) es la capa seleccionada, con un valor comprendido entre 0 y 3.

La detección de los botones se realiza mediante interrupciones GPIO. Cuando se produce un cambio de estado, la rutina de interrupción introduce el número del GPIO en una cola de FreeRTOS. El procesamiento se realiza después en una tarea normal, evitando incluir operaciones complejas dentro de la interrupción.

La tarea espera 50 milisegundos antes de leer los controles. Este pequeño retraso cumple dos funciones: permite detectar combinaciones en las que los dedos no pulsan exactamente al mismo tiempo y reduce los efectos del rebote mecánico de los botones.

El sistema utiliza además una variable de desplazamiento para seleccionar el banco de caracteres. Un desplazamiento de 0 activa las cuatro capas de letras, mientras que un desplazamiento de 4 activa las capas de números y símbolos.

La capa real se obtiene sumando ambos valores:

\begin{equation}
    C_{\text{real}} = C + D
\end{equation}

donde \(D\) puede valer 0 o 4.

Una de las posiciones de la última capa contiene el valor especial \texttt{0xFF}. Al detectar este valor, el programa cambia el desplazamiento entre 0 y 4, alternando así entre los dos bancos de caracteres.

También se ha implementado una combinación formada por los dos botones y el pulsador del joystick. Esta combinación permite alternar entre el modo de escritura y el modo de exploración. En el modo de exploración, los caracteres aparecen en la matriz, pero no se envían al dispositivo conectado.

\subsection{Selección y repetición de caracteres}
\label{subsec:seleccion-caracteres}

La relación entre capas, direcciones y caracteres se almacena en dos matrices de ocho filas y ocho columnas. Una de ellas contiene los códigos HID que deben enviarse, mientras que la segunda contiene la representación textual utilizada por la matriz RGB.

Cuando se detecta una dirección, el programa calcula la capa real y consulta ambas tablas. De esta forma obtiene tanto el código que se enviará mediante Bluetooth como el carácter que debe aparecer en la matriz.

Se ha añadido un mecanismo para evitar que una misma dirección produzca varias pulsaciones por error. Cuando aparece una dirección nueva, esta debe detectarse durante varias iteraciones antes de aceptarse.

En el caso de algunas teclas, como espacio, Intro y borrar, sí se permite la repetición. Si el joystick permanece en la misma dirección durante aproximadamente dos segundos, el programa comienza a repetir la pulsación cada 100 milisegundos. Este comportamiento se ha limitado a las acciones donde mantener una tecla pulsada resulta útil.

El sistema también contempla teclas que necesitan modificadores. Algunos símbolos requieren enviar una combinación con \textit{Shift} o \textit{AltGr}. Para representar estas combinaciones se almacena tanto el código HID de la tecla como el modificador necesario.

\subsection{Visualización en la matriz RGB}
\label{subsec:implementacion-matriz}

La matriz RGB se utiliza para mostrar el carácter seleccionado antes o durante su envío. Cada uno de sus 60 LED representa un píxel dentro de una cuadrícula de seis columnas y diez filas.

Para dibujar los caracteres se ha integrado una fuente gráfica de \(6 \times 10\) píxeles. La fuente contiene 256 glifos y cada uno ocupa diez bytes, uno por cada fila de la matriz. Cada bit indica si un píxel debe estar encendido o apagado.

La función de dibujo recibe el carácter que se quiere mostrar, el color, la capa activa y el estado de mayúsculas. Después calcula la posición del glifo dentro de la fuente y recorre sus diez filas y seis columnas, actualizando el estado de cada LED.

También se ha añadido una conversión para tratar caracteres codificados en UTF-8 que no se encuentran directamente en el rango ASCII. Este tratamiento se utiliza especialmente para mostrar correctamente la letra \textit{ñ} y su versión mayúscula.

Cuando el bloqueo de mayúsculas se encuentra activo, las letras minúsculas se convierten a su equivalente en mayúscula antes de consultar la fuente.

Además del carácter, la primera columna de la matriz muestra un color relacionado con la capa activa. Esto permite reconocer visualmente en qué capa se encuentra el teclado sin necesidad de mostrar información adicional.

La matriz no está pensada para enseñar frases completas. Su función consiste en mostrar un único carácter y ofrecer una confirmación inmediata de la selección realizada.

\subsection{Comunicación Bluetooth HID}
\label{subsec:implementacion-bluetooth}

La comunicación inalámbrica se ha implementado mediante Bluetooth Low Energy utilizando la pila NimBLE. El dispositivo se anuncia con el nombre \enquote{Teclado Unimano} y utiliza el perfil HID para ser reconocido como un teclado.

El módulo Bluetooth registra un servicio HID y define un descriptor estándar de teclado. Los informes enviados están formados por ocho bytes: un byte para modificadores, un byte reservado y seis posiciones para códigos de tecla.

Para generar una pulsación completa se realizan dos envíos. El primero contiene el código de la tecla y, cuando es necesario, su modificador. Tras una espera de 30 milisegundos, se envía un segundo informe vacío para indicar que la tecla ha sido liberada.

Además del servicio HID, se han incorporado los servicios de información del dispositivo y nivel de batería. En la versión actual, el nivel de batería anunciado es un valor fijo, por lo que todavía queda pendiente conectarlo a una medición real de la tensión de la batería.

El emparejamiento utiliza almacenamiento no volátil, permitiendo conservar la información del dispositivo conectado después de un reinicio. Si la conexión se pierde, el teclado vuelve a anunciarse automáticamente.

Una limitación de esta implementación es que algunos símbolos dependen de la distribución de teclado configurada en el sistema operativo. Las combinaciones utilizadas se han planteado para una distribución española, por lo que podrían producir resultados diferentes si el equipo receptor utiliza otra distribución.

\section{Fabricación de la carcasa}
\label{sec:fabricacion-carcasa}

La carcasa se está diseñando mediante FreeCAD y será fabricada utilizando impresión 3D. El diseño debe permitir sujetar el dispositivo con una sola mano, mantener el joystick al alcance del pulgar y colocar los botones en una posición cómoda para los dedos índice y corazón.

También debe reservar espacio para el XIAO ESP32C3, la matriz RGB, la batería, el sistema de alimentación y el cableado interno.

% TODO: Completar después de fabricar la carcasa:
% - Material utilizado para la impresión.
% - Altura de capa, relleno y soportes.
% - Dimensiones finales.
% - Número de versiones impresas.
% - Cambios realizados entre versiones.
% - Fotografías o capturas del modelo de FreeCAD.

\section{Integración del prototipo}
\label{sec:integracion-prototipo}

La integración final consistirá en soldar las conexiones, fijar los componentes dentro de la carcasa y comprobar que los controles pueden utilizarse sin que el cableado limite el cierre del dispositivo.

La matriz deberá quedar visible desde el exterior, mientras que el puerto de carga y el interruptor de encendido deberán permanecer accesibles. La batería tendrá que fijarse de forma que no pueda moverse dentro de la carcasa.

% TODO: Sustituir este texto por el proceso real cuando se complete:
% - Orden seguido durante el montaje.
% - Forma de fijación de cada componente.
% - Organización de los cables.
% - Acceso al USB-C y al interruptor.
% - Resultado final del prototipo.

\section{Problemas encontrados y soluciones}
\label{sec:problemas-soluciones}

Durante el desarrollo del firmware se encontraron varios problemas relacionados con la lectura de entradas y la representación de caracteres.

El primer problema fue la inestabilidad propia de las señales analógicas del joystick. Incluso cuando se encontraba en reposo, los valores del ADC presentaban pequeñas variaciones. Para evitar que estas diferencias se interpretasen como movimientos se definió una zona central con un margen de tolerancia.

También aparecieron dificultades al detectar pulsaciones simultáneas de los botones. Los dedos no accionan ambos controles exactamente en el mismo instante y, además, los botones producen rebotes eléctricos. La solución adoptada fue utilizar interrupciones, una cola de eventos y una espera de 50 milisegundos antes de leer el estado conjunto.

Otro problema fue la repetición involuntaria de caracteres cuando el joystick permanecía en una dirección durante varias lecturas. Para solucionarlo se almacenó la última tecla detectada y se exigió que el joystick volviese al centro antes de aceptar de nuevo la misma letra. La repetición continua quedó habilitada únicamente para espacio, Intro y borrar.

La representación de caracteres como la \textit{ñ} también necesitó un tratamiento específico. La fuente gráfica utilizada emplea una codificación distinta a UTF-8, por lo que fue necesario interpretar las secuencias de varios bytes y convertirlas al índice correspondiente dentro de la fuente.

Finalmente, los símbolos que utilizan \textit{Shift} o \textit{AltGr} exigieron almacenar información adicional junto al código HID. Su comportamiento depende de la distribución de teclado del sistema receptor, por lo que todavía será necesario comprobar todos los símbolos durante la fase de pruebas.

% TODO: Añadir aquí los problemas encontrados durante:
% - La soldadura.
% - El montaje dentro de la carcasa.
% - La alimentación con batería.
% - La autonomía.
% - Las primeras pruebas de uso.



\end{document}
